% math0 manuscript template. Build with: latexmk -pdf paper.tex
\documentclass[11pt]{amsart}
\usepackage[margin=1.1in]{geometry}
\usepackage{amsmath,amssymb,amsthm,mathtools}
\usepackage{hyperref}
\usepackage{xcolor}
\usepackage{listings}
\usepackage{newunicodechar}

\lstdefinelanguage{lean}{
  keywords={theorem,lemma,def,abbrev,example,fun,by,have,show,from,let,in,if,then,else,match,with,
    namespace,end,open,import,section,variable,instance,structure,where,calc,obtain,intro,exact,
    refine,rcases,rintro,simp,omega,decide,norm_num,linarith,nlinarith,positivity,induction,cases},
  sensitive=true, comment=[l]{--}, morecomment=[s]{/-}{-/}, string=[b]"}
\lstset{language=lean, basicstyle=\ttfamily\small, columns=fullflexible, breaklines=true,
  keywordstyle=\color{blue!60!black}\bfseries, commentstyle=\color{gray}, frame=single,
  extendedchars=true, inputencoding=utf8}
% Unicode commonly found in Lean sources.
\newunicodechar{ℕ}{\ensuremath{\mathbb{N}}}
\newunicodechar{ℤ}{\ensuremath{\mathbb{Z}}}
\newunicodechar{ℚ}{\ensuremath{\mathbb{Q}}}
\newunicodechar{ℝ}{\ensuremath{\mathbb{R}}}
\newunicodechar{ℂ}{\ensuremath{\mathbb{C}}}
\newunicodechar{∀}{\ensuremath{\forall}}
\newunicodechar{∃}{\ensuremath{\exists}}
\newunicodechar{→}{\ensuremath{\to}}
\newunicodechar{←}{\ensuremath{\leftarrow}}
\newunicodechar{↔}{\ensuremath{\leftrightarrow}}
\newunicodechar{↦}{\ensuremath{\mapsto}}
\newunicodechar{¬}{\ensuremath{\neg}}
\newunicodechar{∧}{\ensuremath{\wedge}}
\newunicodechar{∨}{\ensuremath{\vee}}
\newunicodechar{≤}{\ensuremath{\le}}
\newunicodechar{≥}{\ensuremath{\ge}}
\newunicodechar{≠}{\ensuremath{\ne}}
\newunicodechar{∈}{\ensuremath{\in}}
\newunicodechar{∉}{\ensuremath{\notin}}
\newunicodechar{⊆}{\ensuremath{\subseteq}}
\newunicodechar{∩}{\ensuremath{\cap}}
\newunicodechar{∪}{\ensuremath{\cup}}
\newunicodechar{∑}{\ensuremath{\sum}}
\newunicodechar{∏}{\ensuremath{\prod}}
\newunicodechar{⟨}{\ensuremath{\langle}}
\newunicodechar{⟩}{\ensuremath{\rangle}}
\newunicodechar{·}{\ensuremath{\cdot}}
\newunicodechar{∘}{\ensuremath{\circ}}
\newunicodechar{⁻}{\ensuremath{^{-}}}
\newunicodechar{¹}{\ensuremath{^{1}}}
\newunicodechar{²}{\ensuremath{^{2}}}
\newunicodechar{₀}{\ensuremath{_{0}}}
\newunicodechar{₁}{\ensuremath{_{1}}}
\newunicodechar{₂}{\ensuremath{_{2}}}
\newunicodechar{α}{\ensuremath{\alpha}}
\newunicodechar{β}{\ensuremath{\beta}}
\newunicodechar{ε}{\ensuremath{\varepsilon}}
\newunicodechar{δ}{\ensuremath{\delta}}
\newunicodechar{π}{\ensuremath{\pi}}
\newunicodechar{σ}{\ensuremath{\sigma}}
\newunicodechar{φ}{\ensuremath{\varphi}}
\newunicodechar{∞}{\ensuremath{\infty}}
\newunicodechar{⌊}{\ensuremath{\lfloor}}
\newunicodechar{⌋}{\ensuremath{\rfloor}}
\newunicodechar{⌈}{\ensuremath{\lceil}}
\newunicodechar{⌉}{\ensuremath{\rceil}}
\newunicodechar{∣}{\ensuremath{\mid}}
\newunicodechar{≡}{\ensuremath{\equiv}}
\newunicodechar{×}{\ensuremath{\times}}
\newunicodechar{⧸}{\ensuremath{/}}
\newunicodechar{‖}{\ensuremath{\|}}
\newunicodechar{√}{\ensuremath{\surd}}
\newunicodechar{∅}{\ensuremath{\emptyset}}
\newunicodechar{⦃}{\ensuremath{\{\!|}}
\newunicodechar{⦄}{\ensuremath{|\!\}}}
\newunicodechar{𝓝}{\ensuremath{\mathcal{N}}}
\newunicodechar{≫}{\ensuremath{\gg}}
\newunicodechar{≪}{\ensuremath{\ll}}

\theoremstyle{plain}
\newtheorem{theorem}{Theorem}
\newtheorem{lemma}[theorem]{Lemma}
\newtheorem{proposition}[theorem]{Proposition}
\newtheorem{corollary}[theorem]{Corollary}
\theoremstyle{definition}
\newtheorem{definition}[theorem]{Definition}
\newtheorem{remark}[theorem]{Remark}

% listings does not handle multi-byte UTF-8 input; map the characters used below.
\lstset{literate={∀}{{$\forall$}}1 {ℕ}{{$\mathbb{N}$}}1 {→}{{$\to$}}1 {≤}{{$\le$}}1
  {·}{{$\cdot$}}1 {⊆}{{$\subseteq$}}1, keepspaces=true, basicstyle=\ttfamily\footnotesize}
\setlength{\emergencystretch}{3em}
\newcommand{\floor}[1]{\left\lfloor #1\right\rfloor}
\newcommand{\Tri}{\mathcal{T}}
\newcommand{\adj}{\sim}
\newcommand{\lean}[1]{\texttt{#1}}

\title[Rademacher's theorem in Lean 4]{Graphs with $\lfloor n^2/4\rfloor+1$ edges contain at least $\lfloor n/2\rfloor$ triangles: a Lean~4 proof of Rademacher's theorem}
\author{math0}
\address{Produced by Claude (an AI model developed by Anthropic) in the math0 project}
\date{October 17, 2026}

\begin{document}
\begin{abstract}
We give a complete, self-contained proof of Rademacher's theorem: every simple graph on $n$
vertices with exactly $\lfloor n^2/4\rfloor+1$ edges contains at least $\lfloor n/2\rfloor$
triangles. The proof is a strong induction over vertex subsets. It derives Mantel's theorem
along the way and splits on the degree sums along the edges of a triangle. The result is
classical. It is the case $t=1$ of Erdős Problem~\#1010, whose general form was proved by
Lovász--Simonovits and by Nikiforov--Khadzhiivanov, and a Lean proof of that general form
already exists. Our contribution is therefore an independent formal proof of a known result. It
is verified in Lean~4 with Mathlib against the statement
\texttt{Erdos1010.erdos\_1010.variants.rademacher} of the formal-conjectures repository and uses
only the standard axioms.
\end{abstract}
\maketitle

\section{Introduction}

Mantel's theorem~\cite{Mantel} says that a graph on $n$ vertices with more than
$\lfloor n^2/4\rfloor$ edges contains a triangle. The bound is attained by the balanced complete
bipartite graph $K_{\lceil n/2\rceil,\lfloor n/2\rfloor}$. In 1941 Rademacher proved, without
publishing it, that one edge beyond this threshold already forces many triangles: every graph on
$n$ vertices with $\lfloor n^2/4\rfloor+1$ edges contains at least $\lfloor n/2\rfloor$ triangles.
The result first appeared in print, with a short proof, in a paper of Erdős~\cite{Erdos55}; see also~\cite{Erdos62} and the account
in~\cite{LiuPikhurko}. The bound is sharp. Add one edge inside the larger part of
$K_{\lceil n/2\rceil,\lfloor n/2\rfloor}$ (with $n\ge 3$). Every triangle of the resulting graph
contains the added edge, and its third vertex lies in the other part, so there are exactly
$\lfloor n/2\rfloor$ triangles.

Erdős~\cite{Erdos62} extended Rademacher's theorem: a graph on $n$ vertices with
$\lfloor n^2/4\rfloor+t$ edges contains at least $t\lfloor n/2\rfloor$ triangles whenever $t<cn$,
for some absolute constant $c>0$. He asked whether the same holds for all
$t<\lfloor n/2\rfloor$. This is Erdős Problem~\#1010 on the website maintained by
Bloom~\cite{EP1010}. The question was answered affirmatively, independently by Lovász and
Simonovits~\cite{LovaszSimonovits} and by Nikiforov and Khadzhiivanov~\cite{NikiforovKh}.
The more general Erdős--Rademacher problem asks for the minimum number of triangles in a graph of
given order and size. It has a large literature. For large $n$ and edge density bounded away
from $1$, it was solved exactly by Liu, Pikhurko and Staden~\cite{LPS}, whose introduction surveys
the history.

\subsection*{Formal statements and prior formalization}
The formal-conjectures repository~\cite{FC} formalizes Problem~\#1010 in Lean~4 as
\texttt{Erdos1010.erdos\_1010}. It also records Rademacher's case $t=1$ as a separate statement,
\texttt{Erdos1010.erdos\_1010.variants.rademacher}, which is the statement proved here. Both are
tagged \texttt{research solved}. The upstream file attaches a \texttt{formal\_proof} link to the
general statement \texttt{erdos\_1010}. The link points to a Lean~4 proof in the repository
\texttt{plby/lean-proofs}~\cite{plby} (file
\nolinkurl{src/latest/ErdosProblems/Erdos1010.lean}, commit \texttt{8822f7d}). That proof covers
every finite vertex type and every $t<\lfloor n/2\rfloor$. Taking $t=1$ gives the present
statement for $n\ge 4$. The remaining cases are trivial (Remark~\ref{rem:small}). The variant
itself carries no \texttt{formal\_proof} attribute upstream, but a formal proof of a statement
that implies it was already available. We therefore do not claim a first formalization.

What this note adds is an independent machine-checked proof of the variant, stated exactly as
upstream. The proof is a short, self-contained double induction and needs about 330 lines of
Lean. According to its header, the existing development uses maximum-cut charge bounds and a
reduction for odd orders. The argument formalized here is elementary, and we make no claim that it
is new as a piece of mathematics. Its first half is the standard inductive proof of Mantel's
theorem.

\section{Statement}

All graphs are finite and simple. For a graph $G$ we write $u\adj v$ when $u$ and $v$ are
adjacent, $e(G)$ for the number of edges, and $t(G)$ for the number of triangles, i.e.\ of
$3$-element vertex sets that span a complete graph.

\begin{theorem}[Rademacher]\label{thm:main}
Let $n\ge 0$ and let $G$ be a simple graph on the vertex set $\{0,\dots,n-1\}$ with
$e(G)=\lfloor n^2/4\rfloor+1$. Then $t(G)\ge\lfloor n/2\rfloor$.
\end{theorem}

The upstream statement is copied below verbatim from the file
\nolinkurl{FormalConjectures/ErdosProblems/1010.lean} of formal-conjectures at commit
\texttt{c77fdd4}. It sits inside \lean{namespace Erdos1010}, with \lean{open SimpleGraph} in
effect:
\begin{lstlisting}[basicstyle=\ttfamily\scriptsize]
open scoped Classical in
/--
Rademacher proved that every graph on $n$ vertices with $\lfloor n^2/4\rfloor+1$ edges contains at
least $\lfloor n/2\rfloor$ triangles.
-/
@[category research solved, AMS 5]
theorem erdos_1010.variants.rademacher : ∀ n : ℕ, ∀ G : SimpleGraph (Fin n),
    G.edgeFinset.card = n ^ 2 / 4 + 1 → n / 2 ≤ (G.cliqueFinset 3).card := by
  sorry
\end{lstlisting}
Here \lean{n \textasciicircum{} 2 / 4} and \lean{n / 2} use natural-number division, so they equal
$\lfloor n^2/4\rfloor$ and $\lfloor n/2\rfloor$. The finset \lean{G.cliqueFinset 3} consists of
the $3$-element vertex sets that are cliques, i.e.\ the unordered triangles. The decidability
instances that \lean{edgeFinset} and \lean{cliqueFinset} need come from \lean{open scoped Classical}.
So the formal statement is exactly Theorem~\ref{thm:main}.

\begin{remark}\label{rem:small}
For $n\le 2$ the hypothesis is impossible, because $e(G)\le\binom n2<\lfloor n^2/4\rfloor+1$. For
$n=3$ it forces $G=K_3$, and then $t(G)=1=\lfloor 3/2\rfloor$. The proof below does not treat
these cases separately. The theorem also holds with ``at least $\lfloor n^2/4\rfloor+1$ edges'',
since deleting edges cannot create triangles. In fact the Lean proof establishes the ``at least''
form for every vertex subset (Theorem~\ref{thm:strong}).
\end{remark}

\section{Proof}

Fix a finite vertex set $V$ and a simple graph $G$ on $V$. All sets below are subsets of $V$.

\begin{definition}\label{def:notation}
For $S\subseteq V$ and $v\in V$ (not necessarily in $S$) let $[v\adj w]\in\{0,1\}$ be the
adjacency indicator, and define
\[
  d_S(v)=\sum_{w\in S}[v\adj w],\qquad D(S)=\sum_{v\in S}d_S(v),\qquad
  \Tri(S)=\{U\subseteq S:\ U \text{ is a triangle of } G\}.
\]
So $d_S(v)=|N(v)\cap S|$ and $D(S)=2e(G[S])$. In Lean these are \lean{dg}, \lean{D} and
\lean{tri}; \lean{tri G S} is \lean{(G.cliqueFinset 3).filter (· ⊆ S)}.
\end{definition}

We prove the following stronger statement. It is the Lean lemma \lean{main}.

\begin{theorem}\label{thm:strong}
For every $S\subseteq V$, if $D(S)\ge 2\bigl(\lfloor |S|^2/4\rfloor+1\bigr)$, then
$|\Tri(S)|\ge\lfloor |S|/2\rfloor$.
\end{theorem}

\subsection{Bookkeeping lemmas}

\begin{lemma}[\lean{dg\_erase}, \lean{D\_erase}]\label{lem:erase}
Let $x\in S$. For every $v\in V$, $d_S(v)=[v\adj x]+d_{S\setminus x}(v)$. Moreover
$D(S)=D(S\setminus x)+2d_S(x)$.
\end{lemma}
\begin{proof}
The first identity splits off the term $w=x$. For the second, split off $v=x$ in $D(S)$ and apply
the first identity to the remaining terms:
\[
  D(S)=d_S(x)+\sum_{v\in S\setminus x}\bigl([x\adj v]+d_{S\setminus x}(v)\bigr)
      =d_S(x)+d_{S\setminus x}(x)+D(S\setminus x).
\]
Here we used the symmetry $[v\adj x]=[x\adj v]$. Finally $d_{S\setminus x}(x)=d_S(x)$, because
$[x\adj x]=0$.
\end{proof}

\begin{lemma}[\lean{D\_erase\_two}]\label{lem:erase2}
If $x,y\in S$ and $x\adj y$, then $D(S)+2=D(S\setminus\{x,y\})+2d_S(x)+2d_S(y)$.
\end{lemma}
\begin{proof}
We have $y\in S\setminus x$ since $x\ne y$. Apply Lemma~\ref{lem:erase} twice:
$D(S)=D(S\setminus\{x,y\})+2d_{S\setminus x}(y)+2d_S(x)$. Also
$d_S(y)=[y\adj x]+d_{S\setminus x}(y)=1+d_{S\setminus x}(y)$.
\end{proof}

\begin{lemma}[\lean{sq\_div\_four}]\label{lem:arith}
For every $M\in\mathbb N$, $\lfloor (M+2)^2/4\rfloor=\lfloor M^2/4\rfloor+M+1$.
\end{lemma}
\begin{proof}
$(M+2)^2=M^2+4(M+1)$.
\end{proof}

\begin{lemma}[\lean{step\_bound}]\label{lem:step}
Let $x,y\in S$ with $x\adj y$ and $d_S(x)+d_S(y)\le|S|$, and suppose
$D(S)\ge2(\lfloor|S|^2/4\rfloor+1)$. Put $S'=S\setminus\{x,y\}$. Then $|S'|=|S|-2$ and
$D(S')\ge 2(\lfloor|S'|^2/4\rfloor+1)$.
\end{lemma}
\begin{proof}
Let $M=|S'|$, so $|S|=M+2$. By Lemmas~\ref{lem:erase2} and~\ref{lem:arith},
\[
  D(S')=D(S)+2-2\bigl(d_S(x)+d_S(y)\bigr)
  \ge 2\bigl(\lfloor M^2/4\rfloor+M+2\bigr)+2-2(M+2)=2\bigl(\lfloor M^2/4\rfloor+1\bigr).\qedhere
\]
\end{proof}

\begin{lemma}[\lean{exists\_common}]\label{lem:common}
If $x,y\in V$ and $d_S(x)+d_S(y)>|S|$, then some $z\in S$ satisfies $x\adj z$ and $y\adj z$.
\end{lemma}
\begin{proof}
The sets $N(x)\cap S$ and $N(y)\cap S$ are subsets of $S$ whose sizes add up to more than $|S|$.
By inclusion--exclusion their intersection is nonempty.
\end{proof}

\subsection{Mantel's theorem}

\begin{lemma}[Mantel; \lean{mantel}]\label{lem:mantel}
If $D(S)\ge2(\lfloor|S|^2/4\rfloor+1)$, then there are $x,y,z\in S$ with $x\adj y$, $x\adj z$
and $y\adj z$.
\end{lemma}
\begin{proof}
We use strong induction on $|S|$. Since $D(S)>0$, $S$ contains an edge $xy$
(\lean{exists\_adj\_of\_D\_pos}). If $d_S(x)+d_S(y)>|S|$, Lemma~\ref{lem:common} gives a common
neighbour $z\in S$. Otherwise Lemma~\ref{lem:step} shows that $S\setminus\{x,y\}$ satisfies the
hypothesis. Since it is smaller than $S$, the induction hypothesis gives a triangle inside it,
and hence inside $S$.
\end{proof}

The three vertices are automatically distinct, because $G$ has no loops. So $\{x,y,z\}\in\Tri(S)$
(\lean{mem\_tri}).

\subsection{The two cases}

Let $N=|S|$ and let $xyz$ be a triangle in $S$.

\begin{lemma}[Case 1; \lean{case\_one}]\label{lem:case1}
Suppose $D(S)\ge2(\lfloor N^2/4\rfloor+1)$ and $d_S(x)+d_S(y)\le N$. Put $S'=S\setminus\{x,y\}$
and suppose Theorem~\ref{thm:strong} holds for $S'$. Then $|\Tri(S)|\ge\lfloor N/2\rfloor$.
\end{lemma}
\begin{proof}
By Lemma~\ref{lem:step}, $S'$ satisfies the hypothesis of Theorem~\ref{thm:strong} and
$|S'|=N-2$, so $|\Tri(S')|\ge\lfloor (N-2)/2\rfloor$. Every triangle in $S'$ is a triangle in $S$
(\lean{tri\_mono}). The triangle $\{x,y,z\}$ lies in $\Tri(S)$ but not in $\Tri(S')$, because
$x\notin S'$. Hence $|\Tri(S)|\ge\lfloor (N-2)/2\rfloor+1=\lfloor N/2\rfloor$. The last equality
uses $N\ge2$, which holds because $S$ contains the triangle $\{x,y,z\}$.
\end{proof}

\begin{lemma}[Case 2; \lean{case\_two}]\label{lem:case2}
Suppose $d_S(x)+d_S(y)>N$, $d_S(y)+d_S(z)>N$ and $d_S(x)+d_S(z)>N$. Then
$|\Tri(S)|\ge\lfloor N/2\rfloor$.
\end{lemma}
No assumption on the number of edges is needed here.
\begin{proof}
Let $R=S\setminus\{x,y,z\}$, so $|R|=N-3$. Each of $x,y,z$ is adjacent to the other two and not
to itself. So three applications of Lemma~\ref{lem:erase} give
\begin{equation}\label{eq:dR}
  d_S(x)=2+d_R(x),\qquad d_S(y)=2+d_R(y),\qquad d_S(z)=2+d_R(z).
\end{equation}
Index the three edges of the triangle by $i\in\{0,1,2\}$ as $e_0=\{x,y\}$, $e_1=\{y,z\}$,
$e_2=\{x,z\}$. In Lean, $e_i=\{p_i,r_i\}$ with $p=(x,y,x)$ and $r=(y,z,z)$. Let
\[
  P=\{(w,i)\in R\times\{0,1,2\}:\ w \text{ is adjacent to both endpoints of } e_i\},
  \qquad f(w,i)=e_i\cup\{w\}.
\]

\emph{Counting $P$.} Fix $w\in R$ and let $k_w=[x\adj w]+[y\adj w]+[z\adj w]\in\{0,1,2,3\}$. The
three edges $e_i$ are the three pairs of vertices from $\{x,y,z\}$. So the number of $i$ with
$(w,i)\in P$ equals $\binom{k_w}{2}$, and $k_w\le\binom{k_w}{2}+1$ for $k_w\le3$. Summing over
$w\in R$ gives
\begin{equation}\label{eq:P}
  d_R(x)+d_R(y)+d_R(z)\le|P|+|R|.
\end{equation}
In Lean this is \lean{sum\_le\_sum} followed by a case split over the eight adjacency patterns of
$w$ with $x,y,z$.

\emph{Injectivity.} Each $f(w,i)$ is a triangle in $S$. Its vertices lie in $S$, $e_i$ is an edge,
and $w$ is adjacent to both ends of $e_i$. The map $f$ is injective on $P$. Indeed, suppose
$f(w,i)=f(w',i')$. Then $w$ is the only element of $f(w,i)$ outside $\{x,y,z\}$, so $w=w'$, and
then $e_i=f(w,i)\cap\{x,y,z\}$ determines $i$. In Lean, $w=w'$ is shown first, and then all nine
pairs $(i,i')$ are checked by comparing which of $x,y,z$ belong to the two sets. Finally
$\{x,y,z\}\notin f(P)$, because every element of $f(P)$ contains a vertex of $R$. Therefore
$\{\{x,y,z\}\}\cup f(P)\subseteq\Tri(S)$ is a disjoint union, and
\[
  |\Tri(S)|\ge|P|+1.
\]

\emph{Conclusion.} By \eqref{eq:P} and \eqref{eq:dR},
\[
  |\Tri(S)|\ge d_R(x)+d_R(y)+d_R(z)-|R|+1=d_S(x)+d_S(y)+d_S(z)-6-(N-3)+1.
\]
Adding the three hypotheses gives $2(d_S(x)+d_S(y)+d_S(z))\ge3N+3$. So
$2|\Tri(S)|\ge 3N+3-2N-4=N-1$, and since $|\Tri(S)|$ is an integer,
$|\Tri(S)|\ge\lceil (N-1)/2\rceil=\lfloor N/2\rfloor$. In Lean this last step is a single
\lean{omega} call over natural numbers, using the facts above.
\end{proof}

\subsection{Proof of Theorems~\ref{thm:strong} and~\ref{thm:main}}

\begin{proof}[Proof of Theorem~\ref{thm:strong}]
We use strong induction on $|S|$ (\lean{main}). Lemma~\ref{lem:mantel} gives a triangle $xyz$ in
$S$. If one of its edges, say $uv$, has $d_S(u)+d_S(v)\le|S|$, apply Lemma~\ref{lem:case1} to
that edge and the third vertex. The required instance of the theorem for $S\setminus\{u,v\}$ is
the induction hypothesis, since $|S\setminus\{u,v\}|<|S|$ (\lean{card\_erase\_two\_lt}). In Lean the
three edges are handled by three calls to \lean{case\_one} with the vertices permuted. Otherwise
all three degree sums exceed $|S|$, and Lemma~\ref{lem:case2} applies.
\end{proof}

No separate base case is needed. When $|S|\le 2$ the hypothesis fails. Inside the induction this
is visible because Lemma~\ref{lem:mantel} produces three distinct vertices of $S$.

\begin{proof}[Proof of Theorem~\ref{thm:main}]
Take $V=\{0,\dots,n-1\}$ and $S=V$. By the handshake lemma (Mathlib's
\lean{sum\_degrees\_eq\_twice\_card\_edges}; Lean lemma \lean{D\_univ}), $D(V)=2e(G)=2(\lfloor
n^2/4\rfloor+1)$. So Theorem~\ref{thm:strong} gives $|\Tri(V)|\ge\lfloor n/2\rfloor$, and
$\Tri(V)$ is the set of all triangles of $G$.
\end{proof}

\section{Machine verification}

\subsection*{Files and declarations}
The proof is in the file
\begin{center}
\nolinkurl{lean/Math0/ErdosProblems/Erdos1010_erdos_1010_variants_rademacher.lean}
\end{center}
of the math0 repository (\url{https://github.com/Foundation-0/math0}), in namespace
\lean{Math0.Erdos1010}. It has 333 lines. The main theorem is stated by reference to the upstream
type:
\begin{lstlisting}
theorem erdos_1010_rademacher : type_of% @Erdos1010.erdos_1010.variants.rademacher := by
  intro n G h
  classical
  have key := main G _ (univ : Finset (Fin n)) rfl (by rw [D_univ, h]; simp)
  simpa [tri] using key

end Math0.Erdos1010

#math0_proves Math0.Erdos1010.erdos_1010_rademacher Erdos1010.erdos_1010.variants.rademacher
\end{lstlisting}
The auxiliary development is stated for an arbitrary \lean{Fintype V} with decidable equality and
decidable adjacency. It is instantiated at \lean{Fin n} with the classical instances of the
upstream statement. The correspondence with the text is: Definition~\ref{def:notation}:
\lean{ind}, \lean{dg}, \lean{D}, \lean{tri}; Lemma~\ref{lem:erase}: \lean{dg\_erase},
\lean{D\_erase}; Lemma~\ref{lem:erase2}: \lean{D\_erase\_two}; Lemma~\ref{lem:arith}:
\lean{sq\_div\_four}; Lemma~\ref{lem:step}: \lean{step\_bound}; Lemma~\ref{lem:common}:
\lean{exists\_common} (with \lean{dg\_eq\_card}); Lemma~\ref{lem:mantel}: \lean{mantel}
(with \lean{exists\_adj\_of\_D\_pos}); Lemmas~\ref{lem:case1} and~\ref{lem:case2}:
\lean{case\_one}, \lean{case\_two} (with \lean{mem\_tri}, \lean{tri\_mono});
Theorem~\ref{thm:strong}: \lean{main}; the handshake step: \lean{D\_univ}. The proof involves no
large finite computation. The only case analyses are the eight adjacency patterns in~\eqref{eq:P}
and the nine pairs $(i,i')$ in the injectivity argument. Both are discharged by \lean{by\_cases},
\lean{fin\_cases} and \lean{simp} and are checked by the kernel like any other proof term. The
file contains no \lean{sorry}, \lean{native\_decide}, \lean{axiom} or \lean{set\_option}.

\subsection*{Audit}
The last command of the file, \lean{\#math0\_proves}, is defined in
\nolinkurl{lean/Math0/Audit.lean}. It fails the build unless the following hold.
\begin{enumerate}
\item The type of \lean{Math0.Erdos1010.erdos\_1010\_rademacher} is syntactically identical to the
type of the upstream declaration \lean{Erdos1010.erdos\_1010.variants.rademacher}. The comparison
ignores metadata annotations only: no unfolding and no definitional equality. The universe
parameters must also agree.
\item The theorem depends on no axioms beyond \lean{propext}, \lean{Classical.choice} and
\lean{Quot.sound}. This is computed with Lean's \lean{collectAxioms}. It rules out
\lean{sorryAx}, the axioms behind \lean{native\_decide}, and custom axioms.
\end{enumerate}
On compilation the file reports:
\begin{lstlisting}[language={},basicstyle=\ttfamily\scriptsize]
math0 audit: 'Math0.Erdos1010.erdos_1010_rademacher' depends on axioms [Classical.choice, Quot.sound, propext]
math0 statement check: 'Math0.Erdos1010.erdos_1010_rademacher' proves 'Erdos1010.erdos_1010.variants.rademacher'
\end{lstlisting}

\subsection*{Toolchain and reproduction}
The toolchain is Lean~4~\cite{Lean4} at version \texttt{v4.33.1} with Mathlib~\cite{mathlib} at version \texttt{v4.33.1}. The upstream statements come
from formal-conjectures at commit \texttt{c77fdd446d23539302aa73bcfae4a9fa962b8382}
(\texttt{c77fdd4}), pinned in \nolinkurl{lean/lakefile.toml}. To reproduce, run
\nolinkurl{scripts/verify.sh} from a checkout of the repository with \texttt{elan} installed. It
fetches the Mathlib cache, builds every result with \texttt{lake build}, checks that the
formal-conjectures and Mathlib package checkouts have no local modifications, and reports the
number of statement checks that passed. To check this file alone, run the following in the
\texttt{lean} directory. It takes well under a minute once the dependencies are built.
\begin{lstlisting}[language={}]
lake env lean Math0/ErdosProblems/Erdos1010_erdos_1010_variants_rademacher.lean
\end{lstlisting}

\begin{thebibliography}{10}

\bibitem{EP1010}
T.~F. Bloom, \emph{Erdős Problem \#1010},
\url{https://www.erdosproblems.com/1010}, accessed 2026-10-17.

\bibitem{Erdos55}
P.~Erdős, \emph{Some theorems on graphs}, Riveon Lematematika \textbf{9} (1955), 13--17.

\bibitem{Erdos62}
P.~Erdős, \emph{On a theorem of {R}ademacher--{T}urán}, Illinois J. Math. \textbf{6} (1962),
122--127.

\bibitem{FC}
M.~Firsching, P.~Lezeau, S.~Mercuri, M.~Z. Horváth, Y.~Dillies, C.~Sönne, E.~Wieser, F.~Zhang,
T.~Hubert, B.~Agüera y Arcas, and P.~Kohli, \emph{Formal Conjectures: An open and evolving
benchmark for verified discovery in mathematics}, arXiv:2605.13171, 2026. Repository:
\url{https://github.com/google-deepmind/formal-conjectures}, commit \texttt{c77fdd4}.

\bibitem{LiuPikhurko}
X.~Liu and O.~Pikhurko, \emph{A note on extremal constructions for the
{E}rdős--{R}ademacher problem}, Combin. Probab. Comput. \textbf{34} (2025), no.~1, 52--62;
arXiv:2311.18753.

\bibitem{LPS}
H.~Liu, O.~Pikhurko, and K.~Staden, \emph{The exact minimum number of triangles in graphs with
given order and size}, Forum Math. Pi \textbf{8} (2020), e8.

\bibitem{LovaszSimonovits}
L.~Lovász and M.~Simonovits, \emph{On the number of complete subgraphs of a graph}, in:
Proceedings of the Fifth British Combinatorial Conference (Aberdeen, 1975), Congressus
Numerantium XV, Utilitas Math., Winnipeg, 1976, pp.~431--441.

\bibitem{Mantel}
W.~Mantel, \emph{Problem 28}, Wiskundige Opgaven \textbf{10} (1907), 60--61.

\bibitem{mathlib}
The mathlib Community, \emph{The {L}ean mathematical library}, in: Proceedings of the 9th ACM
SIGPLAN International Conference on Certified Programs and Proofs (CPP 2020), ACM, 2020,
pp.~367--381.

\bibitem{Lean4}
L.~de~Moura and S.~Ullrich, \emph{The {L}ean 4 theorem prover and programming language}, in:
Automated Deduction -- CADE 28, Lecture Notes in Computer Science, vol.~12699, Springer, 2021,
pp.~625--635.

\bibitem{NikiforovKh}
V.~S. Nikiforov and N.~G. Khadzhiivanov, \emph{Solution of the problem of {P}. {E}rdős on the
number of triangles in graphs with $n$ vertices and $[n^2/4]+l$ edges}, C. R. Acad. Bulgare Sci.
\textbf{34} (1981), no.~7, 969--970.

\bibitem{plby}
\texttt{plby/lean-proofs}, \nolinkurl{src/latest/ErdosProblems/Erdos1010.lean}, theorem
\texttt{erdos\_1010}, commit \texttt{8822f7ddef30fadbd92e1c6ab4ed897af356af5e},
\url{https://github.com/plby/lean-proofs}, accessed 2026-10-17.

\end{thebibliography}

\end{document}
